\section*{Hoofdstuk \ref{ch:g9:heredity-and-traits} --- Erfelijkheid en kenmerken}

\begin{solution}{exo:g9:heredity-and-traits:1}
Een \omterm{def:g9:heredity-and-traits:traits}{kenmerk} is elke beschrijfbare eigenschap van een individu;
\omterm{def:g9:heredity-and-traits:heredity}{erfelijkheid} is het doorgeven van \omterm{def:g9:heredity-and-traits:traits}{kenmerken} van ouders aan
nakomelingen via de voortplanting. De sorteertoets: wat twee
\omterm{def:g8:sexual-reproduction:gametes}{geslachtscellen} niet kunnen dragen, kan geen familie doorgeven.
\end{solution}

\begin{solution}{exo:g9:heredity-and-traits:2}
Bloedgroep: \omterm{def:g9:heredity-and-traits:heredity}{erfelijk}. Gaatje in het \omterm{def:g1:the-five-senses:sense}{oor}: verworven. Sproeten:
erfelijke aanleg, in de zon geleefd (beide werkingen). Vlot
Frans: verworven. Vorm van de oorlel: \omterm{def:g9:heredity-and-traits:heredity}{erfelijk}.
\end{solution}

\begin{solution}{exo:g9:heredity-and-traits:3}
Familiepatroon (keert terug langs de lijnen van de
voortplanting, of volgt levens); van het begin af aanwezig
(vast bij de \omterm{def:g5:flower-to-fruit:steps}{bevruchting}, of daarna opgebouwd); doorgeven
(gaat over op kinderen, of sterft met wie het verwierf).
\end{solution}

\begin{solution}{exo:g9:heredity-and-traits:4}
Omdat getrainde \omterm{def:g3:bones-and-muscles:muscle}{spieren} verworven zijn --- door gebruik
opgebouwd, niet door \omterm{def:g8:sexual-reproduction:gametes}{geslachtscellen} gedragen. Het laat zien
dat wat een leven aan een lichaam toevoegt, niet in de
voortplanting wordt geschreven.
\end{solution}

\begin{solution}{exo:g9:heredity-and-traits:5}
Doorgeven wat hij niet toont: de bepaler van het \omterm{def:g9:heredity-and-traits:traits}{kenmerk}
reisde verborgen door hem heen en kwam een generatie later
boven.
\end{solution}

\begin{solution}{exo:g9:heredity-and-traits:6}
Omdat de \omterm{def:g9:heredity-and-traits:traits}{kenmerken} zelf (een kin, een bloedgroep) niet in twee
losse cellen kunnen meerijden --- wat meerijdt moeten
instructies zijn om ze te bouwen.
\end{solution}

\begin{solution}{exo:g9:heredity-and-traits:7}
Geërfd: de marge in lengte die het postuur suggereert --- de
kleinzoon kreeg de bepalers van de familie. Geleefd: rijkere
voeding vulde de marge verder --- de toevoeging van de
\omterm{def:g6:exploring-our-environment:components}{omgeving} boven op de erfenis.
\end{solution}

\begin{solution}{exo:g9:heredity-and-traits:8}
Elk kind krijgt een andere selectie van de bepalers van elke
ouder --- het schudden deelt bij elke \omterm{def:g5:flower-to-fruit:steps}{bevruchting} verse
combinaties. Eeneiige tweelingen splitsen \emph{na} één
\omterm{def:g5:flower-to-fruit:steps}{bevruchting}: één deling van de kaarten, twee opbouwen --- dus
zijn hun verzamelingen bepalers gelijk.
\end{solution}

\begin{solution}{exo:g9:heredity-and-traits:9}
(a) Volgt de \omterm{def:g6:classification-kinship:tree}{stamboom}, verschijnt op de familieleeftijd
ongeacht ambacht of dorp, denkbaar door \omterm{def:g8:sexual-reproduction:gametes}{geslachtscellen}
gedragen: \omterm{def:g9:heredity-and-traits:heredity}{erfelijk}. (b) Door training opgebouwd, volgt de
sport en niet de ouders, geen \omterm{def:g8:sexual-reproduction:gametes}{geslachtscel} draagt
uithoudingsvermogen: verworven --- op een geërfde marge van
bouw en \omterm{def:g4:heart-and-blood:heart}{hart} (\cref{prop:g7:organs-at-work:training} bouwt wat
de \omterm{def:g9:heredity-and-traits:heredity}{erfelijkheid} alleen schetst).
\end{solution}

\begin{solution}{exo:g9:heredity-and-traits:10}
Nooit de vaardigheid --- synapsen worden in elk leven opnieuw
getraind. Wat wél kan overgaan: de grondstoffen van aanleg ---
gehoor, handigheid, temperament --- een geërfde marge die
muzikale huishoudens vervolgens met oefening, les en voorbeeld
vullen (het tweede kanaal, dat van het samenleven).
\end{solution}

\begin{solution}{exo:g9:heredity-and-traits:11}
Overslaan is \omterm{def:g9:heredity-and-traits:heredity}{erfelijkheid} die zich normaal gedraagt: elke ouder
kan de bepaler voor blauwe \omterm{def:g1:the-five-senses:sense}{ogen} ongetoond dragen en doorgeven;
als twee van zulke verborgen exemplaren in één kind samenkomen,
toont het \omterm{def:g9:heredity-and-traits:traits}{kenmerk} zich. Het patroon bewijst dragen zonder
tonen --- en niets over de eerlijkheid van de familie.
\end{solution}

\begin{solution}{exo:g9:heredity-and-traits:12}
Alles wat de vader doorgeeft rijdt mee in de \omterm{def:g8:reproductive-systems:male}{zaadcel}, en de
\omterm{def:g8:reproductive-systems:male}{zaadcel} is een \omterm{def:g1:animals-around-us:parts}{staart} en een samengepakte kern --- dus moeten
de bepalers in de kern zitten. En omdat de vaderlijke
instructies voor een heel mens in een pakketje van duizendsten
van een millimeter passen, moeten ze op moleculaire kleinheid
geschreven zijn --- een schrift fijner dan elke celstructuur
die de schoolmicroscoop toont.
\end{solution}

\begin{solution}{pb:g9:heredity-and-traits:1}
\textbf{1.} Kapelkin: \omterm{def:g9:heredity-and-traits:heredity}{erfelijk} (familiepatroon, van het begin
af aanwezig). Linkshandigheid: erfelijke aanleg (volgt families
losjes, ligt vroeg vast). Bakkershoest: verworven (volgt het
ambacht). Bloedgroepen: \omterm{def:g9:heredity-and-traits:heredity}{erfelijk} (vast bij de \omterm{def:g5:flower-to-fruit:steps}{bevruchting}, door
\omterm{def:g8:sexual-reproduction:gametes}{geslachtscellen} te dragen). Verweerde gezichten: verworven
(volgen het herdersleven).

\textbf{2.} De toets van het familiepatroon: het \omterm{def:g9:heredity-and-traits:traits}{kenmerk} valt
samen met de lijnen van een beroep, niet met die van de
voortplanting --- erin trouwen verwerft het, wegverhuizen raakt
het kwijt, geen \omterm{def:g8:sexual-reproduction:gametes}{geslachtscel} te bekennen.

\textbf{3.} Dat een bepaler gedragen kan worden zonder getoond
te worden: de middelste generatie gaf de instructies voor de
kin door terwijl ze andere kinnen toonde.

\textbf{4.} Elke ouder droeg naast de bepaler van zijn getoonde
groep een verborgen bepaler voor groep O --- en het kind kreeg
de twee verborgen exemplaren. Ongetoond dragen, in de inkt van
het register zelf.

\textbf{5.} Niet de kin, maar haar bepalers --- instructies om
haar te bouwen --- reizend in de enige voertuigen die de
voortplanting heeft: één \omterm{def:g4:animal-reproduction:sexual}{eicel} en één \omterm{def:g8:reproductive-systems:male}{zaadcel} per generatie,
vier overdrachten in de eeuw.

\textbf{6.} In de selecties: elke \omterm{def:g5:flower-to-fruit:steps}{bevruchting} deelde de
gedeelde bepaler van de kin uit naast een verse selectie
half-om-half van al het overige --- dezelfde kaart, andere
handen.

\textbf{7.} Verworven \omterm{def:g9:heredity-and-traits:traits}{kenmerken} worden nooit doorgegeven.
Reden: ze worden door het leven in de weefsels van het lichaam
opgebouwd, en de \omterm{def:g8:sexual-reproduction:gametes}{geslachtscellen} --- twee cellen, die alleen
bepalers dragen --- hebben geen manier om ze vast te leggen.

\textbf{8.} Eén \omterm{def:g5:flower-to-fruit:steps}{bevruchting}, later gesplitst: één uitdeling
van bepalers, twee keer opgebouwd ---
\cref{rem:g5:human-reproduction:twins}, de eeneiige tweeling,
door de overeenkomsten van het register bevestigd.

\textbf{9.} Geslachtscelkanaal: de kin, de bloedgroepen, de
aanleg voor linkshandigheid. Samenleefkanaal: de hoest (de
lucht van de bakkerij), de verwering (het weer van de heuvels)
--- en, uit het albumhoofdstuk, talen, ambachten en
vaardigheden.

\textbf{10.} Uit het overslaan: de middelste generatie droeg de
kin ongetoond. Uit het bloedregister: ouders die A en B toonden
droegen O ongetoond. Twee keer overtreft het dragen het tonen.

\textbf{11.} Fysiek: ze zitten in de kernen van de
\omterm{def:g8:sexual-reproduction:gametes}{geslachtscellen} --- zoveel bewijst de anatomie van de \omterm{def:g8:reproductive-systems:male}{zaadcel}
al. Waarvan ze gemaakt zijn, en hoe zo'n klein pakketje zo'n
groot boek schrijft, is precies de vraag van de volgende
hoofdstukken.

\textbf{12.} De meeste \omterm{def:g9:heredity-and-traits:traits}{kenmerken} zijn geërfde marges, ingeleefd:
de \omterm{def:g8:sexual-reproduction:gametes}{geslachtscellen} delen de schets uit, het leven tekent haar in.
\end{solution}
